\section{Energy-weighted decay of a contraction power}
\label{sec:energy_weighted_power}

The final power step in the algorithm of Arad, Landau, Vazirani and
Vidick~\cite[lines 1811--1815]{Arad2016RigorousRG} requires control of
excitation energy as well as overlap. The following scalar estimates
retain the power threshold and a lower excitation energy. They do not
assert the repeated trimming argument or the final compressed spectral
extraction.

\begin{lemma}[A geometric-sum bound]
    \label{lem:unit_energy_weighted_power}
    \lean{Real.mul_one_sub_pow_le_inv_add_one}
    \leanok
    Let $0\leq x\leq1$ and let $\tau$ be a nonnegative integer. Then
    \begin{align}
        x(1-x)^\tau\leq\frac1{\tau+1}.
    \end{align}
\end{lemma}
\begin{proof}\leanok
    Every term of $\sum_{j=0}^\tau(1-x)^j$ is at least $(1-x)^\tau$.
    Multiplication by $x$ and the geometric-sum identity give
    \begin{align}
        (\tau+1)x(1-x)^\tau
          &\leq x\sum_{j=0}^\tau(1-x)^j
           =1-(1-x)^{\tau+1}\leq1.
    \end{align}
\end{proof}

\begin{lemma}[Excitation energy under a squared power]
    \label{lem:gap_energy_weighted_power}
    \lean{Real.energy_mul_one_sub_div_pow_le}
    \leanok
    \uses{lem:unit_energy_weighted_power}
    Let $A>0$ and $0<\gamma\leq s\leq A$. If
    $(\tau+1)\gamma\geq A$, where $\tau$ is a nonnegative integer, then
    \begin{align}
        s(1-s/A)^{2\tau}\leq\gamma(1-\gamma/A)^\tau.
    \end{align}
\end{lemma}
\begin{proof}\leanok
    Apply the preceding estimate at $x=s/A$. Its first power gives
    $s(1-s/A)^\tau\leq A/(\tau+1)\leq\gamma$.
    The other power is bounded by $(1-\gamma/A)^\tau$, since
    $0\leq1-s/A\leq1-\gamma/A$. Multiplying these two nonnegative
    inequalities proves the result.
\end{proof}

\begin{lemma}[An energy-weighted matrix order estimate]
    \label{lem:matrix_gap_energy_weighted_power}
    \lean{Matrix.IsHermitian.energy_weighted_power_bounds}
    \leanok
    \uses{lem:gap_energy_weighted_power}
    Let $H$ be a Hermitian complex matrix on a finite index set, and
    let $0\leq H\leq AI$, where $A>0$. Suppose $0<\gamma\leq A$ and
    every spectral value of $H$ is either zero or at least $\gamma$.
    If $A\leq(\tau+1)\gamma$, then
    \begin{align}
        0\leq H(I-H/A)^{2\tau}
          \leq\gamma(1-\gamma/A)^\tau I.
    \end{align}
    Empty matrix index sets and $\tau=0$ are permitted.
\end{lemma}
\begin{proof}\leanok
    Put $g(s)=1-s/A$ and $f(s)=sg(s)^{2\tau}$. The order bounds
    on $H$ put its spectrum in $[0,A]$. The preceding scalar estimate
    bounds $f$ at every nonzero spectral value, and the same inequality
    holds at zero. Continuous functional calculus gives
    $f(H)=H(I-H/A)^{2\tau}$ and transports both nonnegative scalar
    inequalities to the asserted matrix inequalities.
\end{proof}
